\documentclass[10pt,a4paper]{amsart}
\usepackage[margin=19mm]{geometry}
\usepackage{amsmath,amssymb,mathtools}
\usepackage[scr=rsfs,cal=euler]{mathalfa}
\usepackage{xfrac}
\usepackage[unicode,hidelinks,bookmarks=false]{hyperref}
\newcommand{\ie}{i.e.\ }
\newcommand{\cf}{cf.\ }
\newcommand{\resp}{respectively\ }
\newcommand{\Subsection}{\subsection}
\providecommand{\nobreakdash}{\nobreak\hbox{-}}
\setlength{\emergencystretch}{2em}
\multlinegap0pt
\usepackage{amssymb}
\newtheorem{lem}{Lemma}
\newtheorem{thm}{Theorem}
\newcommand{\C}{\mathbb C}
\newcommand{\Centre}{\mathrm Z}
\newcommand{\Gund}{G_\bullet}
\newcommand{\HC}{\mathrm{HC}}
\newcommand{\Irr}{\mathrm{Irr}}
\newcommand{\R}{\mathbb R}
\newcommand{\bfe}{\mathbf{e}}
\newcommand{\bff}{\mathbf{f}}
\newcommand{\bfh}{\mathbf{h}}
\newcommand{\catC}{\mathscr C}
\newcommand{\eev}{{\oline{\mathbf{0}}}}
\newcommand{\g}{\mathfrak}
\newcommand{\gund}{\g g_\bullet}
\newcommand{\hund}{\g h_\bullet}
\newcommand{\lund}{\g l_\bullet}
\newcommand{\oline}{\overline}
\newcommand{\sseq}{\subseteq}
\title{Classification of irreducible real modules of real Lie superalgebras}
\author{Siddhartha Sahi, Hadi Salmasian, Vera Serganova}
\date{Source: arXiv:2604.09404v1 (CC BY 4.0)}
\begin{document}
\maketitle

\noindent\emph{Source and licence.} Classification of irreducible real modules of real Lie superalgebras. Version arXiv:2604.09404v1 (CC BY 4.0), distributed by arXiv under the Creative Commons Attribution 4.0 International licence, http://creativecommons.org/licenses/by/4.0/, as recorded in the arXiv licence metadata for this exact version and shown on the abstract page https://arxiv.org/abs/2604.09404v1. Full licence notice: \texttt{licenses/arxiv-2604.09404-CC-BY-4.0.txt}.\par\smallskip

\noindent\textit{Editorial scope.} Counted unit: the article's thm thm:alternateclam (source0.tex lines 2162--2168). The article's complete original proof of it is delivered with it (source0.tex lines 2169--2212, one \texttt{proof} environment of 44 lines). No mathematics is written, repaired or re-derived: the statement and the proof are the article's own lines, in the article's own notation, and no retained line is substituted or re-flowed.  Retained uncounted as the source-local context the proof needs: lines 1902--1908; lines 1939--1942; lines 2012--2015.  Nothing was omitted from any retained span, so the package carries no omission record.  No retained line contains a citation, so the wrapper carries no bibliography and no bibliography entry.  No retained line contains a figure environment, an image-inclusion command or drawing code, so no image is delivered and the package declares no asset dependency.  Editorial formatting only, in addition: an AMS \texttt{amsart} preamble that restates the article's own declarations and macros used by the retained lines, namely the environments \texttt{lem}, \texttt{thm} on separate counters, together with the standard packages those lines need.  The retained environments print in this document as Lemma 1, Theorem 1 in order of appearance, whereas the article prints the same items with the numbering of its own full document.  The complete native source of the article as distributed in the arXiv e-print for this exact version is bundled unchanged as \texttt{sources/198205/source0.tex}.
\begin{lem}
\label{lem:kcompalpha-alpha}
  Let $\g g$ be a complex semisimple Lie algebra and let $\g g^\R$, $\tau$, $\theta$, $\g k^\R$, and $\g p^\R$  be as above.  Let $\g h^\R$ be a compact $\theta$-stable Cartan subalgebra of $\g g^\R$, with complexification $\g h$. Then  
  %involution of $\g k$ such that $\g k^\tau$ is the compact real form of $\g k$. Let $\g h$ be a Cartan subalgebra  of $\g k$ such that $\tau(\g h)= \g h$.  Then
$\tau(\g g_\alpha)=\g g_{-\alpha}$
 for every  $\alpha\in\Phi_\g g$.
 \end{lem}

\begin{lem}
\label{lem:Z(a)has-cpct-Cartan}
Let $\g g$ be a semisimple Lie algebra with a real form $\g g^\R$, and let $\theta$ be a Cartan involution of $\g g^\R$.  Let $\g h^\R=\g t^\R\oplus\g a^\R$ be a $\theta$-stable Cartan subalgebra of $\g g$, and set $\g l^\R:=\Centre_{\g g^\R}(\g a^\R)$. Then $\theta|_{[\g l^\R,\g l^\R]}$ is a Cartan involution of $[\g l^\R,\g l^\R]$ and $\g h^\R\cap [\g l^\R,\g l^\R]$ is  a compact $\theta$-stable Cartan subalgebra of $[\g l^\R,\g l^\R]$. %that is contained in $\g t^\R$. 
\end{lem}

\begin{lem}
\label{lem:roots-real-on-a}
 Every $\alpha\in\Phi_\g g$ is real-valued on $i\g t^\R\oplus \g a^\R$. 
\end{lem}

\begin{thm}
\label{thm:alternateclam}
Let $h_{\beta_1},\ldots,h_{\beta_N}\in\g h_\eev$ denote the coroots of $\beta_1,\ldots,\beta_N$, so that $\beta_i(h_{\beta_i})=2$ for $1\leq i\leq N$. Then for every $[W]\in\Irr(\catC_\g g)$  with $\g b$-highest weight $\lambda\in\g h_\eev^*$ we have 
 \[
 c_\lambda=(-1)^{\sum_{i=1}^N\lambda(h_{\beta_i})}\lambda(\HC_\g b(D_\lambda)).
 \]
\end{thm}

\begin{proof}
We have $\theta(\lund)=\lund$, and the Cartan decomposition of 
$\lund^\R:=\lund\cap\g g^\R$
 is 
\[
\lund^\R=(\lund^\R\cap\g k^\R)\oplus(\lund^\R\oplus \g p^\R)
.\] By Lemma~\ref{lem:Z(a)has-cpct-Cartan},  $\hund^\R\cap\lund^\R$ is a 
compact $\theta$-stable Cartan subalgebra of $\lund^\R$.  Let $\alpha\in\Phi_{\lund}$ and set
\[
\bfe:=e_\alpha\quad,\quad
\bff:=-\tau(e_\alpha)\quad,\quad
\bfh:=-[e_\alpha,\tau(e_\alpha)],
\] 
where the $e_\alpha\in\g l_\alpha$ are choices of root vectors. 
By Lemma~\ref{lem:kcompalpha-alpha} we have  $\tau(e_\alpha)\in\g g_{-\alpha}$, and since $\g a^\R\sseq \ker(\alpha)$ we have $\theta(\alpha)=\alpha$. Furthermore,  
 $\bfh\in i(\hund^\R\cap \lund^\R)$
because $\tau(\bfh)=-\bfh$. 
Lemma~\ref{lem:roots-real-on-a} implies that $\alpha(\bfh)\in\R$. Hence, by suitably scaling $e_\alpha$ and substituting $\alpha$ by $-\alpha$ if necessary, we can assume that $\bfe,\bff,\bfh$ is a  standard $\g{sl}_2$-triple. In matrix notation, the real form of the spanned $\g{sl}_2$ subalgebra that is induced by $\tau$ corresponds to the antilinear involution $X\mapsto -X^*$, and we have 
\[
\bfe=\begin{bmatrix}
0 & 1\\
0 & 0
\end{bmatrix}\quad,\quad
\bff=\begin{bmatrix}
0 & 0\\1 & 0
\end{bmatrix}\quad,\quad
\bfh=\begin{bmatrix}
1 & 0\\0 & -1
\end{bmatrix}.
\]
The algebraic subgroup of $\Gund$ corresponding to this $\g{sl}_2$ subalgebra of $\gund$ is isomorphic to either  $\mathrm{SL}_2(\C)$ or its quotient by $\{\pm I\}$, 
%This is because the subgroup should act on the adjoint module and so its kernel is contained in \pm I. Note that the Cartan subgroup of this (P)SL_2 is contained in the Cartan subgroup H of \Gund: this follows from (i) the Cartan subalgebras satisfy this containment property, and (ii) the exponential map from a Cartan subalgebra is a group homomorphism, hence it is surjective onto the Cartan subgroup. 
and the antiholomorphic involution of $\Gund$ restricts to the map $X\mapsto (X^*)^{-1}$. 
The simple reflection corresponding to $\alpha$ corresponds to 
\begin{equation}
\label{eq:w00}
w_\alpha=\begin{bmatrix}
0 & 1\\-1 & 0
\end{bmatrix}\in\mathrm{SL}_2(\C),
\end{equation}
and we have $w_\alpha\tau(w_\alpha)=w_\alpha^2=-I$, where $I\in\mathrm{SL}_2(\C)$ is the identity matrix. it follows that $\lambda(w_\alpha\tau(w_\alpha))=(-1)^{\lambda(\bfh)}$. 

It is a well known fact (originally due to Kostant) that the longest element of the Weyl group of a semisimple Lie algebra is equal to the product of reflections associated to the roots in the Kostant cascade. These reflections commute because of strong orthogonality. This fact together with the above $\mathrm{SL}_2(\C)$ calculation implies the assertion of the corollary. 
\end{proof}

\end{document}
